\lecture{3.2}{Confounding and Exchangeability}
\label{lec:3-2}

\subsection{Classical Confounder Definition}

The classical epidemiologic description says that a confounder is a variable associated with both exposure and outcome, and not on the causal pathway from exposure to outcome. For example, baseline disease severity may affect both whether a patient receives a risky cancer treatment and whether the patient dies of cancer. Disease severity can therefore confound the treatment-outcome association.

\begin{figure}[H]
\centering
\begin{tikzpicture}[node distance=2.1cm,>=Latex]
  \node[draw, circle, fill=blue!5] (W) {$W$\\severity};
  \node[draw, circle, fill=green!5, below left=of W] (A) {$A$\\treatment};
  \node[draw, circle, fill=orange!8, below right=of W] (Y) {$Y$\\outcome};
  \draw[->, thick] (W) -- (A);
  \draw[->, thick] (W) -- (Y);
  \draw[->, thick] (A) -- (Y);
\end{tikzpicture}
\caption{A classical confounding structure: $W$ is a common cause of treatment and outcome.}
\end{figure}

This definition is useful but incomplete. In complex problems, variables can be time-varying, affected by prior treatment, measured with error, or unmeasured. A more general definition uses counterfactual independence.

\subsection{Counterfactual Independence Definition}

\begin{definition}[No unmeasured confounding for a point treatment]
For a binary treatment $A$, there is no unmeasured confounding given covariates $W$ if
\[
\Yzero \ci A \given W
\qquad\text{and}\qquad
\Yone \ci A \given W.
\]
Equivalently, for $a\in\{0,1\}$,
\[
\Ya{a}\ci A\given W.
\]
\end{definition}

In words, among people with the same covariates $W$, the treatment actually received carries no additional information about the potential outcomes. Within covariate strata, treated and untreated people are exchangeable with respect to what would happen under either treatment.

\subsection{Randomization}

In an ideal randomized trial, treatment assignment is generated by a random device that does not use patient prognosis. If treatment is randomized unconditionally, then
\[
\Yzero \ci A
\qquad\text{and}\qquad
\Yone \ci A.
\]
This means that the treated and untreated groups have the same distribution of potential outcomes before treatment is assigned, up to random sampling variation.

\begin{assumptionbox}
\textbf{Randomization condition.} For a binary point treatment,
\[
\Ya{a}\ci A,\qquad a=0,1.
\]
Plainly, treatment assignment gives no information about what the outcome would be under either treatment level.
\end{assumptionbox}

\subsection{Conditional Randomization}

In observational studies, unconditional randomization is usually implausible. Treatment decisions often depend on covariates. A weaker assumption is that treatment is randomized within strata of measured covariates:
\[
\Ya{a}\ci A\given W,\qquad a=0,1.
\]

This is plausible when $W$ contains all common causes of treatment and outcome needed to make treatment assignment as-if random. For example, if physicians choose early imaging based on baseline pain, neurological symptoms, age, cancer history, and prior utilization, then those variables must be included in $W$ for conditional randomization to be plausible.

\begin{assumptionbox}
\textbf{Conditional randomization.} For every treatment level $a$,
\[
\Ya{a}\ci A\given W.
\]
This says treatment is as-if randomized among people with the same $W$.
\end{assumptionbox}

\subsection{Ignorability}

The terms \emph{ignorability}, \emph{conditional ignorability}, and \emph{no unmeasured confounding} are often used for the same counterfactual independence condition. The idea is that after conditioning on $W$, the treatment assignment mechanism can be ignored for the purpose of identifying the distribution of $\Ya{a}$.

The word "ignored" should not be misunderstood. We do not ignore treatment assignment in estimation. Rather, the assumption says that the remaining dependence between treatment and potential outcomes has been removed by conditioning on $W$.

\subsection{Exchangeability}

Exchangeability is another name for the same idea, emphasizing comparability of groups.

\begin{definition}[Exchangeability]
Treated and untreated groups are \emph{exchangeable} with respect to potential outcomes, possibly conditional on $W$, if their potential outcome distributions are the same:
\[
\Ya{a}\ci A\given W.
\]
\end{definition}

If exchangeability holds, then within a stratum $W=w$, the observed outcomes among people with $A=a$ can represent the potential outcomes under $a$ for everyone in that stratum.

\subsection{Why Ignorability Cannot Be Empirically Verified}

Ignorability involves potential outcomes that are not jointly observed. Among treated people we observe outcomes under treatment, but not what their outcomes would have been without treatment. Among untreated people we observe outcomes without treatment, but not what their outcomes would have been under treatment.

Therefore, the condition
\[
\Ya{a}\ci A\given W
\]
cannot be tested directly from the observed data. We can check whether measured covariates are balanced across treatment groups. We can check whether some covariates predict treatment. But we cannot check whether all unmeasured causes of the outcome have been controlled.

\begin{keyidea}
Exchangeability is a causal assumption, not an empirical fact discovered by a balance table. Balance diagnostics can reveal problems, but they cannot prove that ignorability holds.
\end{keyidea}
